\subsection{Coding-manual pilot and source adjudication}\label{supp:coding-pilot}
The released author encoding (AU) is the coding of record: the author decided every code against the cited passage, using AI only to locate passages and draft candidates. Two later AI passes (AI-1, AI-2) re-coded the same forty-four settings to test the codebook, not the coder. The supplied provenance identifies AI-1 and AI-2 as the same model, not independent people. AI-1 had seen aggregate counts; AI-2 had seen AU's status values for ten settings. Table~\ref{tab:coding-pilot} therefore describes a non-blind coding-manual pilot and sensitivity analysis. Its bootstrap intervals resample settings despite source clustering and dependent encodings; they are not human-coder reliability bounds. \paragraph{Independent human re-coding} The decision-critical subset (44 access classes and protection flags, nine intervention settings, six transfer relations) was released as a blind packet with a codebook. An independent security PhD researcher who had not seen the manuscript or AU coded it from the sources without AI tools, recording a supporting passage for every code. Agreement with AU, before any discussion, is 44/44 for access class ($\kappa=1$), 9/9 for intervention purpose ($\kappa=1$), 4/4 each for protected-condition adaptation and same-intervention utility in the privacy settings ($\kappa=1$), 38/44 for the protection flag ($\kappa=.61$) and 5/6 for transfer relations ($\kappa=.77$); overall 104/111 decisions. The privacy denominator is identical. Four protection-flag disagreements mark with/without contrasts that AU did not code as protection comparisons: the composition section of the communication-graph study, its active front-running setting, and the two web-agent fingerprinting settings with stealth configurations. Two mark AgentTrace settings that the coder read as lacking an unspoofed baseline. The transfer disagreement is T1 (retained versus extended). None of the seven can create a query/workflow privacy setting with both an adapted attacker and matched utility, because neither the communication-graph study nor the identity-obfuscation and blocking settings report original-task utility. The author then adjudicated each disagreement against the cited passages. The coder's code was adopted in six cases: CG-composition (its Sec.~9.12 ladder re-measures the CG-transform properties and is counted with that setting), both Fayolle fingerprinting settings (identity obfuscation), ATR-main and ATR-budget (no unspoofed baseline), and T1 (extended). AU's code was retained for CG-active, whose with/without contrast concerns an integrity endpoint rather than inference. After adjudication, three settings are identity obfuscation, thirty-four have no protection comparison, and the four-setting privacy denominator and the matched-intervention synthesis are unchanged. The coder's sheets, the agreement output and the adjudication record are in the review artifact.

\begin{table*}[t]\centering\small
\caption{Original, unadjudicated coding pilot on forty-four settings. Pairwise columns are Cohen's $\kappa$; the three-pass column is Fleiss' $\kappa$ with a 2,000-draw setting-bootstrap interval. Full agreement in this pilot does not establish field validity.}\label{tab:coding-pilot}
\begin{tabularx}{\textwidth}{@{}Yrrrrr@{}}\toprule
Field & AU--AI-1 & AU--AI-2 & AI-1--AI-2 & Three-pass $\kappa$ [95\%] & All agree \\\midrule
Access, protection, named transfer, open world (each) & 1.00 & 1.00 & 1.00 & 1.00 & 44/44 \\
Failure handling & .46 & .80 & .63 & .63 [.46, .75] & 28/44 \\
Execution outcome & .15 & .46 & .52 & .37 [.23, .48] & 14/44 \\\bottomrule
\end{tabularx}\end{table*}

Removing the ten AI-2-exposed settings leaves thirty-four rows: AU--AI-2 $\kappa$ is .785 for failure handling and .433 for outcomes. This does not remove the shared-model or non-blind design. The number coded \emph{reported} also depends on applicability. Failure handling is 19/42, 18/43 and 16/43 for AU, AI-1 and AI-2; outcomes are 13/38, 7/38 and 15/43. These are descriptive coding ranges with different denominators, not a common-denominator confidence interval. Neither the original 13/38 nor an AI majority is treated as a settled field-coverage estimate.

\paragraph{Rules that resolve the ambiguity}
First name the execution and criterion. A source-defined success/failure count qualifies as \emph{reported}, with the strength of its criterion recorded separately. Explicit parser/repair outcomes qualify for a probe study, not as original-task utility. Persona-alignment diagnostics assess a behavioral component and are \emph{partial}. Raw action labels or final messages alone do not establish an evaluated task outcome. Reserve \emph{not applicable} for a setting with no new execution population; missing success evidence in an execution study is \emph{not located}.

\paragraph{Fifteen disputed records}
Source inspection changes thirteen outcome codes and retains two, with a per-setting ledger. All four ANT settings become \emph{not located}: the original codes were two partial, one reported and one inapplicable, not four inapplicable. Four cloud-benchmark settings become \emph{reported} because Section IV-A states 4,070 successful and 42 failed runs (98.9\%). This is an inherited collection-level criterion, not four independent success measurements; the source's conflicting capture totals remain unresolved. Three PersonaFingerprint settings become \emph{partial} for consistency/diversity diagnostics. AgentPrint behavior becomes \emph{partial} for capability activation/retention, and HTTP signatures becomes \emph{not located} for an evaluated execution outcome.

The communication-graph active setting remains \emph{reported}: Section 9.5/Table 4 supplies live front-running outcomes over 100 rounds per condition, which both AI passes missed. Its offline capture-transformation setting remains \emph{not applicable} because the transformation does not re-execute tasks. These source checks take precedence over votes. The revised sheet contains 15 reported outcomes among 42 applicable rows; twenty-nine rows retain their prior extraction rather than receiving a new independent coding. No post-adjudication $\kappa$ is presented as validation.

The artifact preserves all original encodings, exact agreement calculations, the exposed-row sensitivity, the clarified manual, source versions and the fifteen-row before/after ledger. General execution-outcome reporting remains distinct from original-task utility under the \emph{same defense}; consequently these revisions do not alter the four-setting privacy/utility denominator in the main text.
